\section{Versuche}

\subsection{Versuch 1: Löslichkeit von Silberhalogeniden}

\subsubsection{Durchführung}
In drei Reagenzgläsern wurden Lösungen von Kaliumchlorid, Kaliumbromid und Kaliumiodid hergestellt. Anschließend wurden zu jeder Lösung einige Tropfen Silbernitratlösung gegeben, wodurch Niederschläge entstanden. Im nächsten Schritt wurde untersucht, ob sich die gebildeten Niederschläge durch Zugabe einer konzentrierten Ammoniumcarbonatlösung wieder lösen lassen. Der Versuch wurde mit frischen Ansätzen wiederholt, wobei anstelle der Ammoniumcarbonatlösung konzentrierte Ammoniaklösung verwendet wurde.

\subsubsection{Beobachtung}
Nach Zugabe der Silbernitratlösung entstanden in allen drei Ansätzen Niederschläge. Dabei zeigte Silberchlorid eine weiße Farbe, Silberbromid erschien hellgelb und Silberiodid zeigte eine deutlich gelbe Färbung.

Bei Zugabe der konzentrierten Ammoniumcarbonatlösung löste sich der Silberchloridniederschlag teilweise bis vollständig auf. Der Silberbromidniederschlag zeigte nur eine geringe Löslichkeit, während Silberiodid praktisch ungelöst blieb.

Bei Zugabe konzentrierter Ammoniaklösung löste sich Silberchlorid vollständig auf. Silberbromid löste sich nur teilweise, während bei Silberiodid keine sichtbare Auflösung beobachtet werden konnte.

\subsubsection{Auswertung}
Die Zugabe von Silbernitrat zu den Halogenidlösungen führt zur Bildung schwerlöslicher Silberhalogenide. Für die Fällungsreaktionen ergeben sich die in den Gleichungen \eqref{eq:agcl_fällung}, \eqref{eq:agbr_fällung} und \eqref{eq:agi_fällung} dargestellten Reaktionen.

\begin{equation}
    \ce{Ag^+_{(aq)} + Cl^-_{(aq)} -> AgCl_{(s)}}
    \label{eq:agcl_fällung}
\end{equation}

\begin{equation}
    \ce{Ag^+_{(aq)} + Br^-_{(aq)} -> AgBr_{(s)}}
    \label{eq:agbr_fällung}
\end{equation}

\begin{equation}
    \ce{Ag^+_{(aq)} + I^-_{(aq)} -> AgI_{(s)}}
    \label{eq:agi_fällung}
\end{equation}

Silberchlorid fällt weiß aus, Silberbromid cremefarben und Silberiodid gelb.

Die Auflösung der Niederschläge in Ammoniak beruht auf der Bildung des stabilen Diamminsilber(I)-Komplexes nach \eqref{eq:diamminsilber}.

\begin{equation}
    \ce{Ag^+_{(aq)} + 2 NH3_{(aq)} <=> [Ag(NH3)2]^+_{(aq)}}
    \label{eq:diamminsilber}
\end{equation}

Durch die Komplexbildung wird die Konzentration freier Silberionen verringert, wodurch sich das Löslichkeitsgleichgewicht der Silberhalogenide nach dem Prinzip von Le Chatelier in Richtung Lösung verschiebt. Silberchlorid besitzt das größte Löslichkeitsprodukt der untersuchten Silberhalogenide und löst sich daher am leichtesten. Silberbromid zeigt eine geringere Löslichkeit, während Silberiodid aufgrund seines sehr kleinen Löslichkeitsproduktes praktisch unlöslich bleibt.

Die beobachteten Löslichkeitsunterschiede spiegeln sich in den Löslichkeitsprodukten der Silberhalogenide wider. Diese betragen bei \SI{25}{\celsius} für Silberchlorid etwa $1{,}7 \cdot 10^{-10}$, für Silberbromid etwa $5{,}0 \cdot 10^{-13}$ und für Silberiodid etwa $8{,}5 \cdot 10^{-17}$. \cite{}

Die abnehmenden Löslichkeitsprodukte in der Reihenfolge Silberchlorid > Silberbromid > Silberiodid spiegeln die zunehmende Stabilität der Feststoffe wider und erklären das beobachtete Löslichkeitsverhalten.

Neben der Bildung schwerlöslicher Silbersalze besitzen die Halogene weitere charakteristische periodische Eigenschaften. Sie zeichnen sich durch eine hohe Elektronegativität aus, die innerhalb der Gruppe von Fluor zu Iod abnimmt. Darüber hinaus bilden alle Halogene bevorzugt einwertige Halogenidionen (\ce{F^-}, \ce{Cl^-}, \ce{Br^-} und \ce{I^-}). Eine weitere wichtige Eigenschaft ist die starke oxidierende Wirkung der elementaren Halogene, deren Stärke ebenfalls innerhalb der Gruppe von Fluor zu Iod abnimmt. Diese Gemeinsamkeiten beruhen auf der identischen Valenzelektronenkonfiguration aller Halogene mit sieben Valenzelektronen.

\subsubsection{Fehlerbetrachtung}
Die beobachtete Löslichkeit der Silberhalogenide kann durch mehrere Faktoren beeinflusst worden sein. Bereits geringe Unterschiede bei der eingesetzten Stoffmenge der Kaliumhalogenide oder bei der Anzahl der zugesetzten Silbernitrattropfen führen zu Konzentrationsabweichungen und damit zu Veränderungen der Niederschlagsbildung.

Darüber hinaus kann die Konzentration der verwendeten Ammoniak- beziehungsweise Ammoniumcarbonatlösung von der Sollkonzentration abgewichen sein. Insbesondere Ammoniaklösungen verlieren bei längerer Lagerung durch Verdunstung Ammoniak, wodurch ihre Komplexierungsfähigkeit abnimmt.

Weiterhin können Verunreinigungen der Reagenzgläser oder eine unvollständige Durchmischung der Lösungen die Beobachtungen beeinflussen.

\subsection{Versuch 2: pH-Werte von Hydroxiden und Oxid-Hydroxiden der dritten Periode}

\subsubsection{Durchführung}
In sieben Reagenzgläsern wurden Lösungen von Natriumhydroxid, Magnesiumhydroxid, Aluminiumhydroxid sowie den entsprechenden Oxid-Hydroxiden von Silicium, Phosphor, Schwefel und Chlor hergestellt. Anschließend wurden die pH-Werte der Lösungen mithilfe von Universalindikatorpapier bestimmt.

\subsubsection{Beobachtung}
Die im Versuch bestimmten pH-Werte sind in Tabelle \ref{tab:ph_3_periode} eingetragen.

\begin{table}[H]
    \centering
    \caption{pH-Werte der Hydroxide und Oxid-Hydroxide der dritten Periode}
    \label{tab:ph_3_periode}
    \begin{tabular}{|l|c|}
        \hline
        \textbf{Substanz} & \textbf{pH-Wert} \\
        \hline
        Natriumhydroxid & 13--14 \\
        Magnesiumhydroxid & 9--10 \\
        Aluminiumhydroxid & 5--7 \\
        Siliciumdioxid & 5--6 \\
        Phosphoroxid-Hydroxid & 2--3 \\
        Schwefeloxid-Hydroxid & 0--1 \\
        Chloroxid-Hydroxid & 0--1 \\
        \hline
    \end{tabular}
\end{table}

\subsubsection{Auswertung}
Die beobachteten pH-Werte ergeben sich aus einem systematischen Übergang von stark basischen zu stark sauren Eigenschaften innerhalb der dritten Periode. Ursache hierfür ist die zunehmende effektive Kernladung bei gleichzeitig abnehmendem metallischem Charakter der Elemente. Dadurch verändert sich die Art der Bindung zwischen Zentralatom und Hydroxid- bzw. Oxidliganden grundlegend.

Im Falle stark metallischer Elemente wie Natrium liegt eine vollständige Dissoziation in Wasser vor, wie in Gleichung \eqref{eq:naoh_dissoziation} dargestellt. Die hohe Konzentration an Hydroxidionen führt zu stark alkalischen pH-Werten.

\begin{equation}
    \ce{NaOH_{(s)} -> Na^+_{(aq)} + OH^-_{(aq)}}
    \label{eq:naoh_dissoziation}
\end{equation}

Bei weniger basischen Metallhydroxiden wie Magnesiumhydroxid liegt ein Löslichkeitsgleichgewicht vor, welches in \eqref{eq:mgoh2_gleichgewicht} beschrieben ist. Die geringe Löslichkeit begrenzt die Hydroxidionenkonzentration und damit den pH-Wert.

\begin{equation}
    \ce{Mg(OH)2_{(s)} <=> Mg^{2+}_{(aq)} + 2OH^-_{(aq)}}
    \label{eq:mgoh2_gleichgewicht}
\end{equation}

Amphoteres Aluminiumhydroxid kann sowohl als Säure als auch als Base reagieren. Die Reaktion mit Hydroxidionen ist in Gleichung \eqref{eq:aloh3_base} dargestellt, während die Reaktion mit Protonen in Gleichung \eqref{eq:aloh3_acid} beschrieben ist.

\begin{equation}
    \ce{Al(OH)3_{(s)} + OH^-_{(aq)} <=> [Al(OH)4]^-_{(aq)}}
    \label{eq:aloh3_base}
\end{equation}

\begin{equation}
    \ce{Al(OH)3_{(s)} + 3H^+_{(aq)} <=> Al^{3+}_{(aq)} + 3H2O_{(l)}}
    \label{eq:aloh3_acid}
\end{equation}

Amphoteres Aluminiumhydroxid stellt den Übergangsbereich zwischen basischem und saurem Verhalten dar.

Weiter rechts in der Periode dominiert zunehmend kovalente Bindung zwischen dem jeweiligen Element und Sauerstoff. Die entsprechenden Oxide reagieren in Wasser nicht mehr als Hydroxide, sondern werden durch Hydratation in Oxosäuren überführt. Dieser Zusammenhang lässt sich allgemein durch \eqref{eq:oxidhydrolyse_allgemein_rev} beschreiben.

\begin{equation}
    \ce{E_xO_y + n H2O <=> H_nE_xO_{y+n}}
    \label{eq:oxidhydrolyse_allgemein_rev}
\end{equation}

Dabei bezeichnet \ce{E_xO_y} das jeweilige Elementoxid, \ce{H2O} Wasser und \ce{H_nE_xO_{y+n}} die gebildete Oxosäure in Lösung. Der Faktor $n$ beschreibt die stöchiometrische Anzahl der beteiligten Wassermoleküle.

Die zunehmende Stabilisierung kovalenter \ce{E-O}-Bindungen führt dabei zu einer kontinuierlichen Verschiebung des Säure-Base-Charakters hin zu stark sauren Lösungen.

Für die nichtmetallischen Oxide bzw. Oxid-Hydroxide dominiert somit zunehmend die Bildung von Oxoniumionen in Wasser. Siliciumdioxid kann formal als schwache Kieselsäure beschrieben werden, wie in \eqref{eq:silicic_acid} gezeigt.

\begin{equation}
    \ce{SiO2_{(s)} + 2H2O_{(l)} <=> H4SiO4_{(aq)}}
    \label{eq:silicic_acid}
\end{equation}

Dies führt zu leicht sauren pH-Werten.

Phosphor-, Schwefel- und Chlorverbindungen bilden in Wasser Oxosäuren, die Protonen an das Lösungsmittel abgeben. Für Phosphorsäure ist dies in Gleichung \eqref{eq:phosphorsaure} dargestellt.

\begin{equation}
    \ce{H3PO4_{(aq)} + H2O_{(l)} <=> H3O^+_{(aq)} + H2PO4^-_{(aq)}}
    \label{eq:phosphorsaure}
\end{equation}

Für Schwefelsäure gilt in der ersten Protolyse nahezu vollständige Dissoziation, wie in Gleichung \eqref{eq:schwefelsaure} beschrieben.

\begin{equation}
    \ce{H2SO4_{(aq)} + H2O_{(l)} -> H3O^+_{(aq)} + HSO4^-_{(aq)}}
    \label{eq:schwefelsaure}
\end{equation}

Chlorige Oxosäuren zeigen noch stärkere Säureeigenschaften, was in Gleichung \eqref{eq:chloresaure} veranschaulicht wird.

\begin{equation}
    \ce{HClO3_{(aq)} + H2O_{(l)} -> H3O^+_{(aq)} + ClO3^-_{(aq)}}
    \label{eq:chloresaure}
\end{equation}

Der kontinuierliche Übergang von \eqref{eq:naoh_dissoziation} bis Gleichung \eqref{eq:chloresaure} zeigt, dass mit zunehmendem Nichtmetallcharakter die Fähigkeit zur Bildung von Hydroxidionen abnimmt und stattdessen die Bildung von Oxoniumionen dominiert.

\subsubsection{Fehlerbetrachtung}
Eine mögliche Fehlerquelle ist, dass die pH-Werte nur mithilfe von Universalindikatorpapier bestimmt wurden, wodurch es zu Ableseungenauigkeiten durch subjektive Farbinterpretation kommt. Zudem ist die Genauigkeit bei extremen pH-Bereichen (sehr sauer/sehr basisch) begrenzt. Ein weiterer Fehlerfaktor ist die teilweise geringe Löslichkeit bzw. unvollständige Reaktion der Hydroxide und Oxide, wodurch keine Gleichgewichtskonzentrationen erreicht wurden. Besonders bei Magnesium- und Aluminiumhydroxid kann dies den gemessenen pH-Wert verfälschen.

\section{Zusammenfassung}
